\section{Conclusiones}

\begin{itemize}
  \item \textbf{La misma página sirve para todo.} R-Tree, vocabulario, postings,
        normas, directorio y listas del IVF reutilizan la cabecera de 32 bytes y el
        \code{Pager} del Proyecto 1. Eso hizo que el desglose de I/O por consulta
        ---bloques del índice frente a bloques de la tabla--- fuera comparable entre
        los cuatro módulos y entre los dos paradigmas multimedia sin instrumentación
        adicional.
  \item \textbf{Cada índice gana en su régimen y pierde fuera de él.} El R-Tree deja de
        servir cuando las cajas se solapan (dimensión alta); el IVF cambia recall por
        bloques y la rodilla aparece con pocas celdas; el pre-filtrado relacional gana
        con filtros selectivos y el post-filtrado con filtros laxos. El planificador
        híbrido existe para elegir el régimen con datos del propio índice.
  \item \textbf{Las representaciones aprendidas dominan a las cuantizadas en calidad,
        no necesariamente en costo.} En el Experimento 7 el embedding denso supera
        ampliamente a BoVW en Precision@$k$, mientras que la extracción SIFT es más
        barata y el índice invertido lee pocas páginas por consulta. En el texto, la
        recuperación léxica sigue siendo la mejor para consultas de ítem conocido y la
        semántica la complementa en consultas parafraseadas; por eso el RAG muestra
        ambas rutas lado a lado y permite fusionarlas.
  \item \textbf{La aplicación de IA no necesitó estructuras fuera del motor.} Corpus,
        embeddings, codebook e imágenes indexadas son tablas e índices consultables con
        el mismo SQL del cliente web, y el modelo de lenguaje solo consume los
        fragmentos que el planificador recuperó.
\end{itemize}
